\section{A free quotient in the compact Pell encoding}
\label{sec:free-quotient}

This section gives two additional equations in ten strictly positive
integer variables. Their exact arithmetic costs and total degrees are
\(83/312\) and \(89/278\). The second improves the earlier \(89/288\)
and \(92/278\) points in cost and degree. The first improves the earlier
\(84/328\) point in both metrics. Other tradeoffs remain relevant.
These are ordinary proofs composed with the retained native/refutation
calibration, not new proof-assistant certificates or global optimality
claims.

\subsection{The change of arithmetic architecture}

The incoming constructions used \(Y=qs\) and
\(c=kY+\eta\). The factor \(q\) then propagated through
\(U=ic^2\), making the strong Pell part expensive in degree.
We instead take
\[
 Y=s,\qquad \lambda=q,
\]
and put \(\lambda\) into the two Pell bases. This retains the
needed divisibility while lowering the degree of \(c\).
The complete denominator is recovered only after exponential projection.
The positive shift \(T+\alpha\) is restored: it supplies a direct
early proof of \(b>c\) when \(Y\) is no longer divisible by \(q\).

\subsection{Complete polynomial specification}

The independent variables are
\[
 J,F,\alpha,i,T,w,s,n,\eta,\zeta\in\mathbb Z_{>0}.
\]
The fixed tuple \((\beta,K_0,M_C,M_F)\) is the complete directly
indexed BFS943 radix67 tuple retained with this component. Its largest
port has 592,884 bits. Every byte and its finite reconstruction are
included; a port name is not an extra unknown. In particular
\(\beta\ge31\) is odd and
\(K_0>0\), \(0<M_C<\beta<M_F<2\beta\).
Set
\[
 h=\beta J,\quad q=h+1,\quad C=q-F-\alpha,\quad Y=s.
\]
Choose either the plain chart or the corrected chart:
\[
\begin{array}{c|ccc}
 &W_0&X&D\\ \hline
 \text{plain}&qw&q^2w&q^7\\
 \text{corrected}&w&qw&6q^7
\end{array}
\]
Both use the following polynomial abbreviations:
\begin{align*}
 k&=\eta+\zeta,&c&=kY+\eta,\\
 A&=qY(X+1)+1,&\Delta&=A^2-1,&H&=4A-5,\\
 v&=Xq^2Y^2,&E&=1+v(v+1)k^2,\\
 R&=[h(q-F)+\alpha](q^2-1)+(M_C+qM_F)J,\\
 B&=(K_0+W_0)C-F,&U&=ic^2,\\
 a&=\Delta U^2E,&F_0&=1+a,&G&=EF_0,\\
 F_s&=1+4aF_0,&Q&=\Delta(F_s-1),\\
 b&=c(T+\alpha)-R,&V&=F_sb-c,\\
 K&=X(3c-2X)+2,&L&=bV.
\end{align*}
In the plain chart let
\[
 E_c=k-R-1-q^4s;
\]
in the corrected chart let
\[
 E_c=6(k-R-1-q^4s)-v(R+1)(R-1)(R+3).
\]
Finally define
\begin{align*}
 M&=hDH,\\
 Z&=hD\{H(n-LB)+K\}-H(DB+hE_c),\\
 P&=Q(Z^2-M^2GV^2)+(M^2G-Z^2).
\end{align*}
The requested single equation is \(P=0\). Every displayed abbreviation
is a fixed polynomial. There are no variable exponents, divisions,
factorials, or extra predicates in the emitted circuits.

\paragraph{Theorem.}
For either chart, ZFC proves
\[
 \bigl(\nexists J,F,\alpha,i,T,w,s,n,\eta,\zeta>0:\ P=0\bigr)
 \quad\longleftrightarrow\quad\operatorname{Con}(\mathrm{ZFC}).
\]

\subsection{Ordered soundness argument}

For proof notation only, write
\[
 t=B/h,\quad W=1+hL,\quad \kappa=K/H,
 \quad \sigma=Y/q^3,\quad z=Z/M.
\]
For the plain chart put \(e=k-R-1\), \(\rho=e/q^7\).
For the corrected chart put
\[
 \bar e=k-R-1-\frac{v(R+1)(R-1)(R+3)}6,
 \qquad\rho=\bar e/q^7.
\]
No integrality of these quotients is initially assumed. Identically,
\[
 P=(Q-1)Z^2-M^2G(QV^2-1),\qquad
 z=n+\kappa+\sigma-Wt-\rho.                       \tag{F1}
\]

\paragraph{Root extraction and early positive transport.}
At a positive zero, the identities
\[
 F_s=(1+2a)^2,\qquad Q=4\Delta^2U^2G,
 \qquad\gcd(E,F_0)=1
\]
and prime valuations give \(E=\tau^2\), \(F_0=f_0^2\),
\(G=g^2\), where \(g=\tau f_0>0\). The retained rational
adjacent-Pell descent gives a positive integer \(y\) such that
\[
 QV^2-(Q-1)y^2=1,\qquad z=\pm gy.
\]
With \(f=1+2a\), \(I=2Ug\), one has
\[
 f^2-\Delta I^2=1,\quad c^2\mid I,\quad Q=(\Delta I)^2.
\]
The signed folding theorem gives \(c=\psi_A(p)\) for positive odd
\(p\). Here \(\psi_A(j)\) denotes the Pell ordinate. Since
\(c\ge3\), \(p\ge3\); therefore
\[
 c\ge4A^2-1>4q^4.
\]
The unconditional mask estimate is
\(R<q^4+(\alpha+2)q^2\). Thus
\[
 b-c>\alpha(c-q^2)-q^4-2q^2
       \ge c-q^4-3q^2>0.
\]
This argument includes \(Y=1\) and does not assume a sign of \(C\).
The auxiliary Pell index is odd; because
\(V>c(f^2-1)>4Q\), its first and third values are excluded. Hence
\[
 V>8Q^2,\qquad b>4g,\qquad gy<\frac{3}{8}bV<\frac L2.
                                                               \tag{F2}
\]

For the plain chart, \(\rho<b/q^7<L/2\). If \(B\le-1\),
then \(Wt<-L\), and (F1)--(F2) contradict \(n>0\).
For the corrected chart the same proof applies whenever
\((R+1)(R-1)(R+3)\ge0\). Its remaining cases are \(R\le-4\)
and \(R=0\). The exact identity
\[
 R=(hC+q\alpha)(q^2-1)+(M_C+qM_F)J
\]
shows that, for \(R=-r\le-4\),
\(C<0\), \(|C|>r/q^3\), and
\(Wt<-(1+w|C|)L\). The adverse correction is less than
\(ws^2r^3/(6q^4)\), whereas
\(L>(F_s-1)r^2\) and \(F_s-1>s^2\), so
\(w|C|L\) strictly exceeds it. For \(R=0\) the correction is
only \(ws^2/(2q^4)\) and the same margin suffices.
Therefore both charts force \(B\ge0\), hence \(C>0\).

The calibrated mask bound is now licensed:
\[
 3q<R<q^4-2.
\]
The normalized strong-rank and signed-index theorems, using
\(c^2\mid I\), \(c>2R\), \(c>2p\),
\(V\equiv-c\pmod f\) and \(V\equiv-R\pmod c\), give
\(p=R\). In particular \(R\) is odd.

\paragraph{Authenticate the first index.}
The first square gives \(k=2\psi_{1+2v}(n_0)\).
Comparison of the two Pell bases gives \(1\le n_0<R\).
Because \(H\equiv-1\pmod q\), clearing (F1) in the plain chart
gives \(q^4\mid e\). The first recurrence gives
\(k\equiv2n_0\pmod{4v}\), and here \(q^4\mid v\).
Consequently \(q^4\mid 2n_0-R-1\); its absolute value is less
than \(R<q^4\), so \(2n_0=R+1\).

The corrected chart has only \(q^3\mid v\). Let
\(r_0=(R+1)/2\). The exact Lucas expansion gives
\[
 2\psi_{1+2v}(m)=2m+8v\binom{m+1}{3}+32v^2Z_m,
 \qquad Z_m\in\mathbb Z_{\ge0}.
\]
Thus \(\bar e\) is integral. Clearing (F1) and cancelling the
literal factor 6 gives
\(q^7\mid\bar e-q^4s\), hence \(q^4\mid\bar e\).
For \(d=2n_0-R-1\) one first obtains \(q^3\mid d\), and then
\[
 6(\bar e-d)=vd\{(2n_0)^2+(2n_0)(R+1)+(R+1)^2-4\}
                    +192v^2Z_{n_0}.
\]
The right side is divisible by \(q^6\). For every prime
\(\ell\mid q\), its consequence
\(v_\ell(\bar e-d)\ge6v_\ell(q)-1\ge4v_\ell(q)\)
also covers \(\ell=2,3\). Hence \(q^4\mid d\), and the
same strict bound forces \(d=0\). Thus both charts give
\[
 n_0=(R+1)/2.                                             \tag{F3}
\]

\paragraph{Project before recovering the denominator.}
Reduction of the literal numerator modulo \(H\) gives
\(H\mid hK\) in the plain chart and \(H\mid6hK\) in the
corrected chart. Only the power of \(q\) was cancelled; no coprimality
of \(H\) and 3 is assumed.

Write \(R=2r+1\), \(r\ge48\). The retained ratio theorem is
uniform for every \(\lambda\ge1\), with
\(A=\lambda Y(X+1)+1\) and first base
\(1+2X\lambda^2Y^2\). Apply it at \(\lambda=q\) and use
(F3) and \(Y<c/k<Y+1\). It gives
\[
 Y>X^r/3,\qquad X^{r+1}<3Y(X+1)<H.
\]
The universal projection therefore yields
\(H\mid d_0h(X-2^R)(2^RX+1)\), with \(d_0=1\) or 6.
The entire product has absolute value below \(H\):
\[
 q^3h|(X-2^R)(2^RX+1)|
 <4\cdot4^rX^6+8\cdot16^rX^5<X^{r+1}.
\]
The last inequality holds for \(X\ge32,r\ge48\); after division
by \(X^{r+1}\), the corner value is
\(2^{-117}+2^{-25}<1\), and both terms decrease.
Consequently \(X=2^R\) and \(\kappa\) is integral.

Now \(q\mid X\) implies \(q=2^m\), \(m\ge5\), and
\(R>3q>6m\), so \(q^6\mid X\) and \(q^8\mid v\).
For the plain chart (F3) gives \(4v\mid e\); for the corrected
chart it gives \(\bar e=32v^2Z_{r_0}\). Thus \(\rho\) is a
nonnegative integer in both cases. Finally (F1) gives
\(Y/q^3-WB/h\in\mathbb Z\). Reduction modulo \(q^3\)
and \(h\) proves
\[
 q^3\mid Y,\qquad h\mid B,\qquad B/h\ge0.
\]
The retained parity, exact half-binomial rounding and fresh native-source
reconstruction now give a syntactic ZFC refutation.

\subsection{Constructive converse and exact degree}

A refutation supplies canonical outer data \(q^3\mid Y\),
\(X=2^R\), \(R\equiv3\pmod4\), \(R>3q\), and
\(B/h\ge0\). Set \(s=Y\), and take \(w=X/q^2\) or
\(w=X/q\) for the two charts. Use fresh bases at \(\lambda=q\),
and form \(c=\psi_A(R)\),
\(k=2\psi_{1+2v}((R+1)/2)\),
\(\eta=c-kY\), \(\zeta=k-\eta\).
The uniform ratio gives their strict positivity. A Pell matrix period
modulo \(c^2\tau\) supplies the seed \((f_0,I_0)\), and
\(i=I_0/(c^2\tau)>0\). Double the seed as above. Positive
auxiliary reconstruction gives \(T_*>\alpha\), so
\(T=T_*-\alpha>0\). Set
\[
 n=gy+(1+hL)t+\rho-\kappa-Y/q^3.
\]
This is positive because
\(g\ge f_0>(A-1)I_0>2Xc\), whereas
\(\kappa<Xc\) and \(Y/q^3\le Y<Xc\).
Its carrier equals \(gy\), proving \(P=0\).

For \(q_0=\beta J\), \(k_0=\eta+\zeta\), the complete leading
forms are
\begin{align*}
 [P_{\rm plain}]_{312}
   &=-2^{10}q_0^{140}w^{34}s^{78}i^{14}k_0^{44}(q_0-F)^2,\\
 [P_{\rm corrected}]_{278}
   &=-9\cdot2^{12}q_0^{106}w^{34}s^{78}i^{14}k_0^{44}(q_0-F)^2.
\end{align*}
Both have 135 terms. The highest part of \(b\) comes from \(-R\),
not from \(c(T+\alpha)\); retaining this change is essential.
The independent checker proves 28 exact cut identities per circuit,
covers every gate, verifies these entire leading forms, authenticates
all coefficient bytes, and rejects twenty deliberate actual corruptions.
Normal and optimized Python give identical mathematical receipts.
The complete unbounded proof and retained ordinary dependencies are
provided in the component's \texttt{PROOF.md} and source directory.
